\documentclass[12pt]{article}
\usepackage{notes-project}

\begin{document}

\begin{center}
  {\LARGE\bfseries Chapter 1: Complex Numbers and Euler's Formula}\\[0.25em]
  {\large Mathematical Methods}\\[0.15em]
\end{center}

\section{Complex Numbers} 

\begin{dfn}[Imaginary Unit]{imaginary-number}
An \emph{imaginary unit} is defined as following:
\[
  i^2 = -1
\]
\end{dfn}

\begin{dfn}[Complex Number]{complex-number}
The type of \emph{complex numbers} is defined as following:
\[
  \CC \defeq
  \left\{ a + b \cdot i \,\middle|\, a,b : \mathbb{R} \right\}
\]
An element \(z : \CC\) is a \emph{complex number}
\par\smallskip\noindent{\footnotesize\texttt{LaTeX: \detokenize{\CC}}}
\end{dfn}

For \(z : \CC\), write \(z = a + b \cdot i\), where \(a,b : \mathbb{R}\)

\begin{dfn}[Real Part]{real-part}
The real part function is defined as following:
\[
  \RealPart : \CC \to \mathbb{R}
  \defeq z \map a
\]
\par\smallskip\noindent{\footnotesize\texttt{LaTeX: \detokenize{\RealPart(z)}}}
\end{dfn}

\begin{dfn}[Imaginary Part]{imaginary-part}
The imaginary part function is defined as following:
\[
  \ImagPart : \CC \to \mathbb{R}
  \defeq z \map b
\]
\par\smallskip\noindent{\footnotesize\texttt{LaTeX: \detokenize{\ImagPart(z)}}}
\end{dfn}

\begin{dfn}[Complex Conjugate]{complex-conjugate}
The complex conjugate function is defined as following:
\[
  \Conj : \CC \to \CC
  \defeq z \map a - b \cdot i
\]
\par\smallskip\noindent{\footnotesize\texttt{LaTeX: \detokenize{\Conj(z)}}}
\end{dfn}

\begin{conv}[Complex Conjugate Notation]{complex-conjugate-notation}
For \(z : \CC\), we write
\[
  \conj{z} \defeq \Conj(z)
\]
where \(\Conj\) is the \xref[complex conjugate function]{def:complex-conjugate}
\par\smallskip\noindent{\footnotesize\texttt{LaTeX: \detokenize{\conj{z}}}}
\end{conv}

\section{Arithmetic Operations}

Given two complex numbers $z_1 = a + bi$ and $z_2 = c + di$:

\begin{prop}[Addition:]{addition}
\begin{align}
z_1 + z_2 = (a + c) + (b + d)i
\end{align}
\end{prop}

\begin{prop}[Subtraction:]{subtraction}
\begin{align}
z_1 - z_2 = (a - c) + (b - d)i
\end{align}
\end{prop}

\begin{prop}[Multiplication:]{multiplication}
\begin{align}
z_1 \cdot z_2 &= (a + bi)(c + di) \nonumber \\
&= ac + adi + bci + bdi^2 \nonumber \\
&= ac + adi + bci - bd \nonumber \\
&= (ac - bd) + (ad + bc)i
\end{align}
\end{prop}

\begin{prop}[Division:]{division}
\begin{align}
\frac{z_1}{z_2} = \frac{z_1 \cdot \conj{z_2}}{z_2 \cdot \conj{z_2}} = \frac{(a + bi)(c - di)}{(c + di)(c - di)} = \frac{(ac + bd) + (bc - ad)i}{c^2 + d^2}
\end{align}
\end{prop}

\section{Euler's Formula}

\begin{rem}[Taylor's Expansion]{trigonometric-taylor-series}
Recall from Section 2.3 (Taylor's Expansion), the Taylor series for $\cos\theta$ and $\sin\theta$ around $\theta = 0$:
\begin{align}
\cos\theta &= 1 - \frac{\theta^2}{2!} + \frac{\theta^4}{4!} - \frac{\theta^6}{6!} + \ldots = \sum_{n=0}^{\infty} \frac{(-1)^n \theta^{2n}}{(2n)!} \\
\sin\theta &= \theta - \frac{\theta^3}{3!} + \frac{\theta^5}{5!} - \frac{\theta^7}{7!} + \ldots = \sum_{n=0}^{\infty} \frac{(-1)^n \theta^{2n+1}}{(2n+1)!}
\end{align}
\end{rem}

\begin{dfn}[Definition of $e$]{e}
The natural number $e$ is defined as:
\begin{align}
e = \lim_{n \to \infty} \left(1 + \frac{1}{n}\right)^n = \sum_{n=0}^{\infty} \frac{1}{n!} = 1 + 1 + \frac{1}{2!} + \frac{1}{3!} + \ldots \approx 2.71828
\end{align}
\end{dfn}

\begin{prop}[Taylor series for $e^x$]{exponential-taylor-series}
From this, we can derive the Taylor series for $e^x$ by replacing $1$ with $x$ in each term:
\begin{align}
e^x = \sum_{n=0}^{\infty} \frac{x^n}{n!} = 1 + x + \frac{x^2}{2!} + \frac{x^3}{3!} + \frac{x^4}{4!} + \ldots
\end{align}
\end{prop}

\begin{thm}[Euler's Formula]{eulers-formula}
This is \textbf{Euler's Formula}:
\begin{align}
e^{i\theta} = \cos\theta + i\sin\theta
\end{align}
\end{thm}

\begin{proof}
Substitute $x = i\theta$:
\begin{align}
e^{i\theta} &= 1 + i\theta + \frac{(i\theta)^2}{2!} + \frac{(i\theta)^3}{3!} + \frac{(i\theta)^4}{4!} + \frac{(i\theta)^5}{5!} + \ldots \nonumber \\
&= 1 + i\theta - \frac{\theta^2}{2!} - \frac{i\theta^3}{3!} + \frac{\theta^4}{4!} + \frac{i\theta^5}{5!} - \ldots
\end{align}

Separating real and imaginary parts:
\begin{align}
e^{i\theta} &= \left(1 - \frac{\theta^2}{2!} + \frac{\theta^4}{4!} - \ldots\right) + i\left(\theta - \frac{\theta^3}{3!} + \frac{\theta^5}{5!} - \ldots\right) \nonumber \\
&= \cos\theta + i\sin\theta
\end{align}
\end{proof}

\section{Polar Form of Complex Numbers}

From Section 1.2 (Angles), we know that any point $(x, y)$ in Cartesian coordinates can be expressed in polar coordinates $(r, \theta)$ where:
\begin{align}
x = r\cos\theta, \quad y = r\sin\theta, \quad r = \sqrt{x^2 + y^2}, \quad \theta = \arctan\left(\frac{y}{x}\right)
\end{align}

\begin{dfn}[Polar Form of Complex Numbers]{polar-form}
Therefore, a complex number $z = x + iy$ can be written as:
\begin{align}
z &= r\cos\theta + ir\sin\theta \nonumber \\
&= r(\cos\theta + i\sin\theta) \nonumber \\
&= re^{i\theta}
\end{align}

Where $r = |z|$ is the \textbf{modulus} and $\theta = \arg(z)$ is the \textbf{argument} of the complex number. 
\end{dfn}

\begin{rem}{polar-form-quadrant}
Similar to the transform of a point from Cartesian to Polar coordinates, we must be careful about the quadrant when determining $\theta$ from $\arctan\left(\frac{y}{x}\right)$. 
\end{rem}

\section{Multiplication in Polar Form}

\begin{prop}[Multiplication in Polar Form]{polar-multiplication}
Given two complex numbers $z_1 = r_1 e^{i\theta_1}$ and $z_2 = r_2 e^{i\theta_2}$:
\begin{align}
z_1 \cdot z_2 = r_1 e^{i\theta_1} \cdot r_2 e^{i\theta_2} = r_1 r_2 e^{i(\theta_1 + \theta_2)}
\end{align}

Multiplication in polar form:
\begin{itemize}
    \item Moduli multiply: $|z_1 z_2| = r_1 \cdot r_2$
    \item Arguments add: $\arg(z_1 z_2) = \theta_1 + \theta_2$
\end{itemize}
\end{prop}

\begin{prop}{complex-powers}
For $z^n$ where $z = re^{i\theta}$:
\begin{align}
z^n = r^n e^{in\theta}
\end{align}
\end{prop}

\section{Nth Root}

To find $w = z^{1/n}$, we need to solve $w^n = z$.

\begin{prop}[Nth Root of Complex Numbers]{nth-roots}
This gives exactly $n$ distinct roots:
\begin{align}
z^{1/n} = r^{1/n} e^{i\frac{\theta + 2\pi k}{n}}, \quad k = 0, 1, 2, \ldots, n-1
\end{align}
\end{prop}

\begin{proof}
Let $z = re^{i\theta}$ and $w = \rho e^{i\phi}$. Then:
\begin{align}
w^n = \rho^n e^{in\phi} = re^{i\theta}
\end{align}

Comparing moduli and arguments:
\begin{align}
\rho^n = r \quad &\Rightarrow \quad \rho = r^{1/n} \text{ (unique positive real)} \\
n\phi = \theta \quad &\Rightarrow \quad \phi = \frac{\theta}{n}
\end{align}

However, since angles are periodic with period $2\pi$, we can write $z = re^{i(\theta + 2\pi k)}$ for any integer $k$. Therefore:
\begin{align}
\phi = \frac{\theta + 2\pi k}{n}, \quad k = 0, 1, 2, \ldots, n-1
\end{align}
\end{proof}

\begin{rem}{distinct-nth-roots}
Note: When $k = n$, we get $\frac{\theta + 2\pi n}{n} = \frac{\theta}{n} + 2\pi$, which is the same angle as $k = 0$. So there are exactly $n$ distinct nth roots.
\end{rem}

\end{document}
